\documentclass[11pt,xcolor={dvipsnames},usepdftitle=false]{beamer}
\usepackage{presentation}
\hypersetup{pdfpagemode=UseNone,hidelinks,pdfdisplaydoctitle=true,%
%
pdftitle={Self driving car engineering notebook for Udacity nanodegree}}

% Enter path to PDF file with figures:
%\newcommand{\pdf}{figures.pdf}
\AtBeginSubsection[]
{
  \begin{frame}
    \frametitle{Subsection}
    \centering
    \insertsubsection
    %\usebeamerfont{section title}\insertsectionhead
  \end{frame}
}
\AtBeginSection[]
{
  \begin{frame}
    \frametitle{Section}
    \centering
    \usebeamerfont{section title}\insertsectionhead
  \end{frame}
}

\begin{document}

% Enter title:
\title{Self Driving Car Notebook}

\information%
%
% Enter URL to research paper (can be commented out):
%[https://github.com/pmichaillat/latex-presentation]
%
{Author: TJ}
%
{UDACITY Nanodegree}

%\frame{\titlepage}
\maketitle
\tableofcontents

\section{Introduction}


\begin{frame}
  \frametitle{}


\end{frame}

% %%%%%%%%%%%%%%%%%%%%%%%%%%%%%%%%%%%%%%%%%%%%%%%%%%%%%%%%%%%%%%%%%%% %
% SECTION     MACHINE LEARNING
% %%%%%%%%%%%%%%%%%%%%%%%%%%%%%%%%%%%%%%%%%%%%%%%%%%%%%%%%%%%%%%%%%%% %
\section{Machine Learning}

% %%%%%%%%%%%%%%%%%%%%%%%%%%%%%%%%%%%%%%%%%%%%%%%%%%%%%%%%%%%%%%%%%%% %
% SUBSECTION  MACHINE LEARNING WORKFLOW
% %%%%%%%%%%%%%%%%%%%%%%%%%%%%%%%%%%%%%%%%%%%%%%%%%%%%%%%%%%%%%%%%%%% %
\subsection{Machine Learning Workflow}

% =================================================================== %
% FRAME       ML METRICS
% =================================================================== %
\begin{frame}
  \frametitle{ML Metrics 1/ - }
  \begin{table}
    \begin{tabular}{| c | c |}
      \hline
      \small {\textbf{Metric}} & {\textbf{Definition}}  \\\hline
      \small True Positive (TP)   & Contains item and predicts item   \\\hline
      \small True Negative (TN)   & Does not contain item and predicts no item   \\\hline
      \small False Positive (FP)  & Does not contain item and predicts item   \\\hline
      \small False Negative (FN)  & Contains item and predicts item   \\\hline
    \end{tabular}
  \end{table}

\end{frame}

% =================================================================== %
% FRAME       ML METRICS - CLASSIFICATION
% =================================================================== %
\begin{frame}
  \frametitle{ML Metrics 2/ - Classification}
  \begin{table}
    \begin{tabular}{| c | c | c |}
      \hline
      \small {\textbf{Metric}} & {\textbf{Equation}} & {\textbf{Definition}}  \\\hline
      \small Precision   & $\frac{TP}{TP+FP} $ & Error rate of algorithm when it's predicting an item \\\hline
      \small Recall   & $\frac{TP}{TP+FN} $ & Error rate when the picture is correct  \\\hline
      \small Accuracy   & $\frac{TP+TN}{TP+FN+FP+TN} $ & Total error rate\\\hline
    \end{tabular}
  \end{table}

\end{frame}

% =================================================================== %
% FRAME       ML METRICS - IOU
% =================================================================== %
\begin{frame}
  \frametitle{ML Metrics 3/ - Intersection over Union (IOU)}

  \begin{itemize}
    \small\item Create bounding boxes for prediction and observation (or truth)
    \small\item Ratio of intersection of bounding boxes to union of bounding boxes
    \small\item 50\% ratio considered threshold for true positive
  \end{itemize}

\end{frame}

% TODO : Summarize example

% =================================================================== %
% FRAME       Data Acquisition and Visualization
% =================================================================== %
\begin{frame}
  \frametitle{Data Acquisition and Visualization}

  \begin{itemize}
    \small\item Most time spent at this stage
    \small\item $[$TERM$]$ Domain Gap - an alg trained on one specific dataset may not perform well on another
  \end{itemize}

\end{frame}

% =================================================================== %
% FRAME       Exploratory Data Analysis
% =================================================================== %
\begin{frame}
  \frametitle{Exploratory Data Analysis (EDA)}

  \begin{itemize}
    \small\item Should spend a significant amount of time visually examining dataset 
    \small\item $[$TERM$]$ Domain Shift - Conditions that affect ML alg ability to analyze
      \begin{itemize}
        \small\item Examples: weather, light, sensor change/processing, environment
      \end{itemize}
    \small\item EDA can help tune data analysis for domain shift
    \small\item Vision problem can require looking at 1000s of images 
  \end{itemize}

\end{frame}

% =================================================================== %
% FRAME       Cross Validation
% =================================================================== %
\begin{frame}
  \frametitle{Cross Validation}

  \begin{itemize}
    \small\item Ensures alg will perform well in any environment
    \small\item Evaluates the generalization ability of the model
    \small\item Concepts
      \begin{itemize}
        \footnotesize\item Overfitting
          \begin{itemize}
            \footnotesize\item Model does not generalize
            \footnotesize\item Differentiate training vs test datasets
          \end{itemize}
      \end{itemize}
      \begin{itemize}
        \footnotesize\item Bias-variance tradeoff
          \begin{itemize}
            \footnotesize\item Hard to create a balanced model
            \footnotesize\item $TestError = Var + Bias + \epsilon$
            \footnotesize\item Variance - sensitivity to training data, bias - quality of fit, test error - error rate on test datasit 
            \footnotesize\item Reducing variance will increase bias
          \end{itemize}
      \end{itemize}
      \begin{itemize}
        \footnotesize\item Cross Validation
          \begin{itemize}
            \footnotesize\item Evaluates how well a model generalizes
            \footnotesize\item Be careful of data leakage
          \end{itemize}
      \end{itemize}
  \end{itemize}

\end{frame}

% =================================================================== %
% FRAME       TF Record
% =================================================================== %
\begin{frame}
  \frametitle{TF Record}

  \begin{itemize}
    \small\item Tensorflow custom data format
    \small\item Records in waymo are tf record
    \small\item Use protofile to designate data fields - protocol files
    \small\item Use TF function to create your own custom tf records
  \end{itemize}

\end{frame}

% =================================================================== %
% FRAME       Exercise - 003 - TFRecord;
% =================================================================== %
\begin{frame}
  \frametitle{Exercise - 003 - TFRecord

  \begin{itemize}
    \small\item 
  \end{itemize}

\end{frame}



% %%%%%%%%%%%%%%%%%%%%%%%%%%%%%%%%%%%%%%%%%%%%%%%%%%%%%%%%%%%%%%%%%%% %
% SECTION     EXTRAS
% %%%%%%%%%%%%%%%%%%%%%%%%%%%%%%%%%%%%%%%%%%%%%%%%%%%%%%%%%%%%%%%%%%% %
\section{Extras}
\subsection{Modeling}

% =================================================================== %
% FRAME       VEHICLE MODELING
% =================================================================== %
\begin{frame}
  \frametitle{Vehicle modeling}
  \begin{itemize}
      \small\item Kinetic model
          \begin{itemize}
              \small\item Simplification of dynamic model
              \small\item More tractible
          \end{itemize}
      \small\item Dynamic model
          \begin{itemize}
              \small\item Aim to encompass dynamics of model
              \small\item More robust
          \end{itemize}
  \end{itemize}

\end{frame}

% =================================================================== %
% FRAME       STATE
% =================================================================== %
\begin{frame}
  \frametitle{State}
  \begin{itemize}
      \small\item Position, orientation, velocity
      \small\item In this example
        \begin{itemize}
          \small\item $[$x,y,psi,v$]$
        \end{itemize}
  \end{itemize}

\end{frame}

% =================================================================== %
% FRAME       BUILDING A KINEMATIC MODEL
% =================================================================== %
\begin{frame}
  \frametitle{Building a Kinematic Model}
  \begin{itemize}
    \small\item Actuators - Control inputs 
      \begin{itemize}
        \small\item eg steering wheel, pedals
        \small\item in example, these are delta for turning and a for acceleration
      \end{itemize}
    %\small\item Formulas
  \end{itemize}

\end{frame}

\begin{frame}
  \frametitle{Following a Trajectory}
  \begin{itemize}
    \small\item With self-driving car
      \begin{itemize}
        \small\item Perception - localization - path planning - control
        \small\item Usually use third degree polynomial for planning
        \small\item Localization compares model to map with SDC
        \small\item Path planning uses map, veh loc, environmental model
      \end{itemize}
  \end{itemize}

\end{frame}

\begin{frame}
  \frametitle{Error 1/}
  \begin{itemize}
    \small\item Must measured and expected values
    \small\item State becomes $[$x,y,psi,v,cte,epsi$]$
    \small\item CTE - Crosstrack error
      \begin{itemize}
        \small\item Low complexity heuristic
        \small\item Difference between intended route and current track
        \small\item For next cte $$cte_{t+1} = f(x_t) - y_t +  (v_t * sin(e\psi_t)*dt) $$
        \small\item Where $cte_{t+1} = f(x_t) - y_t$ is curr cte and rest is error due to veh mvmt
      \end{itemize}
  \end{itemize}

\end{frame}

\begin{frame}
  \frametitle{Error 2/}
  \begin{itemize}
    \small\item Orientation error
      \begin{itemize}
        \small\item Difference between desired and curr orientation
        \small\item $$e\psi_{t+1} = e\psi_t + \frac{v_t}{L_f}*\delta_t*dt$$
        \small\item Becomes
        \small\item $$e\psi_{t+1} = e\psi_t - \psi des_t + (\frac{v_t}{L_f}*\delta_t*dt)$$
        \small\item Where $e\psi_{t+1} = e\psi_t - \psi des_t$ is the curr orient error and rest is new error caused by movement
      \end{itemize}
  \end{itemize}

\end{frame}


%\begin{frame}
%\frametitle{Generic slide}
%\begin{itemize}
%\item Consectetur adipiscing elit
%\item Sed do eiusmod tempor incididunt
%\item Quis nostrud exercitation ullamco laboris nisi ut aliquip ex ea commodo consequat
%\item Duis aute irure dolor in reprehenderit in voluptate velit esse cillum dolore eu fugiat nulla pariatur
%\item Ut labore et dolore magna aliqua
%\end{itemize}
%\end{frame}


\end{document}
